% Problema 3, parte 2: campi da incollare nella piattaforma (uno per sezione)

% ===== CAMPO: 1. Result and scope =====
Complete proof that D\_B(T\_k) ≥ k\^{}2 − k for all k ≥ 1 with the
explicit extremal γ\_k = (k−1,k−1,k−2,\ldots,2,1,1) (Theorem A), resting
on a new exact cell-motion lemma for B (Lemma 1) that is reusable for
the other cells; exhaustive exact verification that D\_B(T\_k) = k\^{}2
− k for k ≤ 9. The general upper bound D\_B(T\_k) ≤ k\^{}2 − k is not
proved.


% ===== CAMPO: 2. Proof =====
\subparagraph{Conventions}\label{conventions}

Draw a partition \(\lambda=(\lambda_1\ge\dots\ge\lambda_s)\) as a
diagram whose \textbf{column} \(i\) (\(1\le i\le s\)) consists of the
cells \((i,h)\), \(1\le h\le \lambda_i\) (\(h\) = height). The
\emph{diagonal} of a cell is \(d(i,h)=i+h-1\). Diagonal \(d\) contains
the \(d\) \emph{slots} \((p,\,d+1-p)\), \(p=1,\dots,d\); a set of cells
is the diagram of a partition iff the column heights are weakly
decreasing. \(\delta_k\) is the diagram whose diagonals \(1,\dots,k\)
are full and all others empty. Throughout \(n=T_k\).

\subparagraph{\texorpdfstring{Lemma 1 (exact cell rule for
\(B\))}{Lemma 1 (exact cell rule for B)}}\label{lemma-1-exact-cell-rule-for-b}

Let \(\lambda\) have \(s\) columns. Then \(B(\lambda)\) is the diagram
obtained from \(\lambda\) by moving every cell as follows:

\begin{enumerate}
\def\labelenumi{\arabic{enumi}.}
\tightlist
\item
  \((i,1)\mapsto(1,i)\) for \(1\le i\le s\);
\item
  \((i,h)\mapsto(i+1,h-1)\) if \(2\le h\le s+1\);
\item
  \((i,h)\mapsto(i,h-1)\) if \(h\ge s+2\).
\end{enumerate}

In particular cells of type 1--2 stay on their diagonal (slot
\(p\mapsto p+1\), and slot \(d\mapsto 1\): a cyclic rotation of the
diagonal), while cells of type 3 move from diagonal \(d\) to diagonal
\(d-1\) (same column).

\emph{Proof.} By definition the columns of \(B(\lambda)\) are the
multiset \(\{s\}\cup\{\lambda_i-1:\lambda_i\ge 2\}\). Consider first the
``unsorted'' diagram \(U\): column 1 of height \(s\), column \(i+1\) of
height \(\lambda_i-1\) (\(1\le i\le s\)). Rules 1 and 2 applied to all
cells with \(h\le s+1\) produce exactly the cells of \(U\) of height
\(\le s+1\) (the new column, and the old columns shifted right and
lowered), and rule 2 without the restriction \(h\le s+1\) would produce
all of \(U\). Let
\(q=\#\{i:\lambda_i-1\ge s+1\}=\#\{i:\lambda_i\ge s+2\}\); because
\(\lambda\) is sorted, these are \(i=1,\dots,q\), so in \(U\) exactly
the columns \(2,\dots,q+1\) have height \(>s+1\) (column 1 has height
\(s\), columns \(i+1\) with \(i>q\) have height \(\le s\)). Now compare
\(U\) with the diagram \(V\) obtained from \(U\) by moving every cell of
height \(\ge s+1\) in columns \(2,\dots,q+1\) one column to the left.
Column heights of \(V\): column 1 has the \(s\) cells of the new column
plus the cells of heights \(s+1,\dots,\lambda_1-1\) received from column
2, hence height \(\lambda_1-1\); column \(i\) for \(2\le i\le q\) keeps
heights \(1,\dots,s\) (from old column \(i-1\), whose height
\(\lambda_{i-1}-1\ge s+1\)), loses heights \(s+1,\dots,\lambda_{i-1}-1\)
and receives heights \(s+1,\dots,\lambda_i-1\) from column \(i+1\),
hence height \(\lambda_i-1\); column \(q+1\) keeps heights \(1,\dots,s\)
and loses the rest, hence height \(s\); columns \(i+1>q+1\) are
unchanged with height \(\lambda_i-1\le s\). So the columns of \(V\) are
\(\lambda_1-1\ge\dots\ge\lambda_q-1\ge s\ge \lambda_{q+1}-1\ge\dots\),
i.e.~\(V\) is the sorted diagram of \(B(\lambda)\) (columns of height 0
at the end are discarded). Finally, the cells moved in passing from
\(U\) to \(V\) are exactly the images under rule 2 of the cells
\((i,h)\) of \(\lambda\) with \(h-1\ge s+1\), i.e.~\(h\ge s+2\);
composing ``\((i,h)\mapsto(i+1,h-1)\) then one column left'' gives rule
3. \(\square\)

\subparagraph{\texorpdfstring{Lemma 2 (convergence to \(\delta_k\); used
only for context, not for the
bound)}{Lemma 2 (convergence to \textbackslash delta\_k; used only for context, not for the bound)}}\label{lemma-2-convergence-to-delta_k-used-only-for-context-not-for-the-bound}

Let \(W(\lambda)=\sum_{\text{cells}} d(i,h)\). By Lemma 1,
\(W(B(\lambda))=W(\lambda)-\#\{(i,h)\in\lambda: h\ge s+2\}\le W(\lambda)\).
Since diagonal \(d\) has only \(d\) slots, among all sets of \(T_k\)
cells with at most \(d\) cells on diagonal \(d\) the minimum of \(W\) is
attained only by \(\delta_k\) (fill diagonals \(1,\dots,k\)); hence
\(W(\lambda)\ge W(\delta_k)\) with equality iff \(\lambda=\delta_k\).
(Combined with the fact that a partition \(\ne\delta_k\) of \(T_k\)
cannot be periodic with all steps of cost 0 --- which is C1 material and
not needed below --- this gives convergence; we do not use it.)

\subparagraph{\texorpdfstring{Theorem A (lower bound). For every
\(k\ge1\), \(D_B(T_k)\ge k^2-k\), attained
by}{Theorem A (lower bound). For every k\textbackslash ge1, D\_B(T\_k)\textbackslash ge k\^{}2-k, attained by}}\label{theorem-a-lower-bound.-for-every-kge1-d_bt_kge-k2-k-attained-by}

\[\gamma_k=(k-1,\,k-1,\,k-2,\,k-3,\dots,2,\,1,\,1)\quad(k\ge2),\qquad \gamma_1=(1).\]
(\(\gamma_k\) is \(\delta_k\) with the part \(k\) replaced by the two
parts \(k-1\) and \(1\); \(|\gamma_k|=T_k\).)

\emph{Proof.} For \(k=1\) the only partition is \(\delta_1\) and
\(D_B(1)=0=k^2-k\). Let \(k\ge2\). For \(p_H\in\{1,\dots,k\}\) and
\(p_X\in\{1,\dots,k+1\}\) let \(S(p_H,p_X)\) be the set of cells of
\(\delta_k\) with the cell \((p_H,\,k+1-p_H)\) removed (a \emph{hole} in
slot \(p_H\) of diagonal \(k\)) and the cell \((p_X,\,k+2-p_X)\) added
(an \emph{extra cell} in slot \(p_X\) of diagonal \(k+1\)). Its column
heights are \(k+1-p\) for \(p\notin\{p_H,p_X\}\), \(k-p_H\) for
\(p=p_H\), \(k+2-p_X\) for \(p=p_X\) (and column \(k+1\) exists iff
\(p_X=k+1\)). It is a partition diagram iff \(p_X\ne p_H\) and
\(p_X\ne p_H+1\) (if \(p_X=p_H\) the cell \((p_H,k+2-p_H)\) would sit
above a missing cell; if \(p_X=p_H+1\) column \(p_X\) would be taller
than column \(p_H\); otherwise the heights are weakly decreasing because
column \(p_X\) has the same height as column \(p_X-1\) and column
\(p_H\) has the same height as column \(p_H+1\)). Note
\(\gamma_k=S(1,k+1)\): columns \(k-1,k-1,k-2,\dots,2,1,1\).

Number of columns of \(S(p_H,p_X)\): \(s=k+1\) if \(p_X=k+1\); \(s=k-1\)
if \(p_H=k\) (then \(p_X\le k\)); \(s=k\) otherwise. Heights: every
column has height \(\le k\) except column \(p_X\), of height
\(k+2-p_X\le k+1\), with equality iff \(p_X=1\). Hence a cell of height
\(\ge s+2\) exists only when \(s=k-1\) and \(p_X=1\), i.e.~only in
\(S(k,1)\), and then it is the single cell \((1,k+1)\).

\emph{Case \(S(k,1)\).} By Lemma 1 with \(s=k-1\): the bottom row
\((i,1)\), \(i\le k-1\), becomes column 1 with heights \(1,\dots,k-1\);
the cells \((i,h)\) with \(2\le h\le k\) (all remaining cells of
\(\delta_k\) minus the hole, since the hole is \((k,1)\) and column
\(k\) is empty) go to \((i+1,h-1)\), giving columns \(2,\dots,k\) of
heights \(k-1,\dots,1\); the cell \((1,k+1)\) drops to \((1,k)\),
completing column 1 to height \(k\). So \(B(S(k,1))=\delta_k\).

\emph{Other cases.} No cell has height \(\ge s+2\), so by Lemma 1 every
cell rotates one slot along its diagonal; the same is true of the empty
slot (hole) on diagonal \(k\), because the cells of diagonal \(k\)
permute cyclically among its slots. Hence \(B(S(p_H,p_X))=S(p_H',p_X')\)
with \(p_H'\equiv p_H+1\pmod k\), \(p_X'\equiv p_X+1\pmod{k+1}\)
(representatives in \(\{1..k\}\), \(\{1..k+1\}\)). (Lemma 1 guarantees
the result is a partition diagram, so the validity condition is
automatically preserved.)

Therefore, starting from \(\gamma_k=S(1,k+1)\), as long as no step has
hit \(S(k,1)\) we have \(B^t(\gamma_k)=S(p_H(t),p_X(t))\) with
\(p_H(t)\equiv 1+t\pmod k\), \(p_X(t)\equiv k+1+t\equiv t\pmod{k+1}\).
None of these equals \(\delta_k\) (they have a hole). The first
\(t\ge0\) with \(S(p_H(t),p_X(t))=S(k,1)\) satisfies
\(t\equiv k-1\pmod k\) and \(t\equiv1\pmod{k+1}\). Writing \(t=k-1+ak\):
\(k-1+ak\equiv -2-a\equiv 1\pmod{k+1}\), so
\(a\equiv-3\equiv k-2\pmod{k+1}\), and the least \(a\ge0\) is \(a=k-2\)
(valid since \(k\ge2\)). Thus \(t_0=k-1+(k-2)k=k^2-k-1\),
\(B^{t}(\gamma_k)\ne\delta_k\) for \(0\le t\le k^2-k-1\), and
\(B^{k^2-k}(\gamma_k)=B(S(k,1))=\delta_k\). Since \(\delta_k\) is a
fixed point and (by Lemma 2, or by C1) it is the only cyclic partition
of \(T_k\), \(d_B(\gamma_k)=k^2-k\). \(\square\)

(Consistency check with the statement's example and data:
\(\gamma_3=(2,2,1,1)\), \(d_B=6\); the exhaustive computation below
gives \(D_B(T_k)=k^2-k\) and confirms \(\gamma_k\) among the extremals
for \(k\le9\).)

\subparagraph{Theorem B (upper bound) --- NOT PROVED
HERE}\label{theorem-b-upper-bound-not-proved-here}

Claim: \(d_B(\lambda)\le k^2-k\) for every partition \(\lambda\) of
\(T_k\). This is the theorem of Igusa (1985) / Etienne (1991) (CITED,
proof not reproduced). What is established:

\begin{itemize}
\tightlist
\item
  \textbf{Exhaustive verification for \(k\le9\)} (exact integer
  arithmetic, all \(p(T_k)\) partitions, 0.5 s total):
  \(D_B(T_k)=k^2-k\) for \(k=1,\dots,9\). This proves the cell's
  statement only for \(k\le 9\).
\item
  Structural facts from Lemma 1 available for the next attempt: (i)
  \(W\) strictly decreases exactly at steps where some pile has size
  \(\ge s+2\), and \(\gamma_k\) has \(W(\gamma_k)-W(\delta_k)=1\), so
  the difficulty is bounding runs of ``free'' steps, not the number of
  costly steps; (ii) empirically (exhaustive, \(k\le8\)) the number of
  piles lies in \(\{k-1,k,k+1\}\) for all \(t\ge T_{k-1}-(k-2)\), which
  suggests the strategy: bound the time until the configuration becomes
  ``near-staircase'', then analyse the near-staircase dynamics as in
  Theorem A.
\end{itemize}

\subparagraph{Conclusion of this
attempt}\label{conclusion-of-this-attempt}

\(D_B(T_k)\ge k^2-k\) for all \(k\ge1\) (Theorem A, complete), with
explicit extremal \(\gamma_k\); \(D_B(T_k)=k^2-k\) for \(k\le9\) by
exhaustive computation; the general inequality \(D_B(T_k)\le k^2-k\)
remains to be proved.


% ===== CAMPO: 3. Verification: instructions, dependencies, timings =====
Python 3 standard library. Scripts (re-run by the orchestrator, see
\texttt{runs/p3\_c2/verifica/}): - code\_1 (python, rigor
\texttt{exact}): Exhaustive exact computation of D\_B(T\_k) and of all
extremal partitions for k ≤ 9 (all partitions of T\_k, orbit followed to
δ\_k). Finite set covered: every partition of T\_k for k=1..9 (up to
89134 partitions). Wall-clock 0.48 s. Confirms D\_B(T\_k)=k\^{}2−k for
k≤9; does not prove the general upper bound. - code\_2 (python, rigor
\texttt{exact}): Exploration only: along every orbit for k ≤ 8, records
the last time the pile count leaves \{k−1,k,k+1\} and the first time
diagonals 1..k−1 are full; used to guide the (unfinished) upper-bound
strategy. Exact integers, all partitions of T\_k for k=3..8,
\textasciitilde1 min.


% ===== CAMPO: 4. Sources and contribution =====
\begin{itemize}
\tightlist
\item
  B. Hopkins, 30 Years of Bulgarian Solitaire, College Math. J. 43
  (2012) 135--140 (read in full; p.137 statement that Igusa proved
  γ\_k=(k−1,k−1,k−2,\ldots,2,1,1) is at maximal distance k(k−1)) ---
  CITED for context only
\item
  K. Igusa, Solution of the Bulgarian solitaire conjecture, Math. Mag.
  58 (1985) 259--271 --- NOT read (paywalled); statement cited via
  Hopkins/Drensky
\item
  G. Etienne, Tableaux de Young et solitaire bulgare, J. Combin. Theory
  Ser. A 58 (1991) 181--197 --- NOT read
\item
  J. R. Griggs, C.-C. Ho, The cycling of partitions and compositions
  under repeated shifts, Adv. Appl. Math. 21 (1998) 205--227 --- NOT
  read
\item
  arXiv:1503.00885 (Drensky) and arXiv:2607.17194 (Meštrović) --- read;
  both only cite Igusa/Etienne for the k(k−1) bound and name Toom's
  extremal τ=γ\_k
\item
  M. Jonsson, Processes on Integer Partitions and Their Limit Shapes,
  PhD thesis, Mälardalen Univ. 2017 (DiVA diva2:1082060) --- read the
  relevant pages; cites Etienne for the game-tree height k\^{}2−k
\item
  Position with respect to the literature: The listed arXiv abstracts
  (math/0401385, 1703.07102: random variants; 1101.1546: Toom's
  convergence proof; 2208.14496: orbit growth; 1503.00885 and
  2607.17194: surveys) do not prove the bound. I fetched and read
  (pdftotext) Drensky 1503.00885, Meštrović 2607.17194, Hopkins ``30
  years of Bulgarian solitaire'' (College Math. J. 43 (2012) 135--140),
  Hopkins--Jones (EJC 13 (2006) R80), N. Pham's honors thesis and M.
  Jonsson's PhD thesis (DiVA 1082060). All state, CITED: Knuth
  conjectured and Igusa (Math. Mag. 58 (1985) 259--271) and Etienne
  (JCTA 58 (1991) 181--197) proved that for n=T\_k the maximal number of
  moves is k(k−1), attained by γ\_k=(k−1,k−1,k−2,\ldots,2,1,1) (Hopkins
  2012, p.~137: ``Igusa {[}15{]} shows that the partition γ\_k \ldots{}
  is at distance k(k−1) from τ\_k and that this distance is maximal'').
  None of the read sources reproduces the proof. My approach follows the
  classical ``cards on diagonals'' idea mentioned by Hopkins (p.~137)
  but makes it exact (Lemma 1) and uses it to prove the lower bound in
  full; the upper bound argument of Igusa/Etienne is not reproduced.
\end{itemize}


% ===== CAMPO: 5. Limits and unresolved parts =====
\begin{itemize}
\tightlist
\item
  Upper bound D\_B(T\_k) ≤ k\^{}2 − k for all k: not proved (only
  verified exhaustively for k ≤ 9 and cited from Igusa 1985 / Etienne
  1991, whose proofs I could not access). This is the next blocker.
\item
  Theorem A uses that δ\_k is the only cyclic partition of T\_k to
  conclude d\_B(γ\_k) = k\^{}2−k rather than merely
  B\textsuperscript{\{k}2−k\}(γ\_k)=δ\_k; this is the C1 statement,
  sketched via the potential W in Lemma 2 but not written out in full
  here (the strict-decrease/non-periodicity step is omitted).
\item
  Lemma 1's proof treats ties (columns of equal height) implicitly: the
  sorted diagram is determined by the multiset of column heights, so
  ties do not affect the cell set, but the reader may want this stated
  explicitly.
\end{itemize}

